\subsection{Overview}

Denitrifying woodchip bioreactors---buried beds or walls of carbon-rich media that intercept subsurface (tile) drainage and support microbial denitrification---have become one of the best-characterised edge-of-field practices for reducing nitrate export from artificially drained cropland. The evidence base spans controlled microcosm and column studies, pilot- and full-scale field installations, and global syntheses, and it consistently establishes that woodchip bioreactors remove nitrate reliably across a wide range of climates and designs. Consolidating full-size systems worldwide, \textcite{Christianson2021} report a mean nitrate-N load reduction of $40 \pm 26\%$ (median 46\%; $n=15$) for bed-style bioreactors and nitrate-N removal rates on the order of $7.2 \pm 9.6$ g N m$^{-3}$ d$^{-1}$ (median 5.1 g N m$^{-3}$ d$^{-1}$; $n=27$), and the practice has been assessed as removing between 20 and 40\% of annual nitrate-N loss in the Midwest. The largest single compilation to date, covering 36 site-years at ten full-size Illinois bioreactors, found a median annual removal of 49 kg N per bioreactor per year (4.1 kg N ha$^{-1}$ y$^{-1}$), a median edge-of-field N removal efficiency of 17\%, and removal in all years, reinforcing the dependability of the practice while pointing to room for design optimisation \parencite{Christianson2026}.

Four broad lines of evidence organise what is known. First, field-scale monitoring quantifies nitrate concentration reductions, load reductions, and volumetric removal rates and shows their dependence on hydraulic residence time (HRT), temperature, flow, and nitrate loading. Second, mechanistic and microbial studies establish denitrification as the dominant removal pathway, sustained over years by woodchip carbon and enhanced denitrifier populations. Third, an active design literature refines sizing, hydraulics, and performance-enhancement strategies (alternative carbon sources, P-sorbing media, pumping, electrical stimulation, and mineral amendments). Fourth, work on secondary effects documents by-products---nitrous oxide, methane, sulfate reduction, dissolved organic carbon (DOC), and phosphorus dynamics---and generally concludes that these ``pollution-swapping'' risks are real but small at landscape scale. The remainder of this section synthesises these lines of evidence, with the recurring limitation that additional study is still needed for phosphorus, pesticides, and greenhouse gases \parencite{Mitchell2023}.

\subsection{Field-scale nitrate removal: rates, efficiencies, and load reductions}

Field installations consistently reduce nitrate, but the magnitude spans a wide, well-explained range. In a controlled Illinois field experiment, bioreactors reduced nitrate loads by an average of 63\%, from a minimum of 20\% at 2-h HRT to a maximum of 98\% at 8-h HRT, with an average removal rate of 11.6 g NO$_3$-N m$^{-3}$ d$^{-1}$ (range 5--30) \parencite{Bell2015}. Under cold-climate conditions in Quebec, a three-year evaluation reported effluent nitrate-N reduced by 99\% and total-N by 72--73\% relative to influent, a 99\% reduction in flow-weighted nitrate-N load, and median removal rates of 6.84 g NO$_3$-N m$^{-3}$ d$^{-1}$ (by volume) and 2.99 g NO$_3$-N m$^{-2}$ d$^{-1}$ (by area); nitrate-N represented 84\% of total-N in the influent, non-compliance events fell by 96\%, and roughly 55\% of drainage passed through the bioreactors each year \parencite{Husk2017}. Across four South Dakota field bioreactors, nitrate-N concentration reductions ranged from 7\% to 100\% with removal rates of 5.2--26.5 g N m$^{-3}$ d$^{-1}$ and load reductions of 39--89\% \parencite{Thapa2023}. A New Zealand system lowered outflow far below inflow, with removal efficiency of 31\% at 2.8 d HRT rising to 93\% at 6.4 d HRT \parencite{Rivas2020}, and a Missouri bioreactor typically produced effluent below 10 mg NO$_3$-N L$^{-1}$, cutting spring influent means of 17 mg L$^{-1}$ to 5 mg L$^{-1}$ and post-fertilisation means of 69 mg L$^{-1}$ to 21 mg L$^{-1}$ \parencite{Aide2016}.

Long-term and multi-site datasets place these numbers in context. In the Illinois ten-bioreactor compilation, systems removed a median 51\% of the N routed into them at a median annual rate of 4.6 g N m$^{-3}$ d$^{-1}$; standardised designs achieved a median flow-weighted HRT of 2.8 h (near the 3.0-h design target), and only annual N-loss removal differed significantly between standardised and hydraulic-maximised designs (4.7 versus 3.2 kg N ha$^{-1}$; $p=0.031$) \parencite{Christianson2026}. An early nine-year Iowa denitrification-wall study found annual nitrate loss of 24.5 kg NO$_3$-N ha$^{-1}$ under the wall versus 54.5 kg ha$^{-1}$ under conventional drainage, with no significant trend over time in either treatment \parencite{Moorman2010}. Not all field systems meet performance targets: a Virginia bioreactor at the low end of N inputs cumulatively removed only an estimated 9.5\% of influent NO$_3$-N (0.56 $\pm$ 0.25 g m$^{-3}$ d$^{-1}$) in its second year---greater in the growing (20.4\%) than non-growing (5.4\%) season---below the NRCS 25--45\% target, with removal rates from $-2.25$ to 8.10 g m$^{-3}$ d$^{-1}$ and net export on 17\% of paired-sample days \parencite{Bock2018}. Similarly, a paired Illinois system removed only 6--28\% of edge-of-field nitrate (1--55\% for individual units), equivalent to 1.8--5.6 kg N ha$^{-1}$ (52--161 kg N field$^{-1}$) annually, and met the federal 20\% edge-of-field target in only one of three years \parencite{Maxwell2022}. Controlled column and laboratory studies confirm the achievable ceiling: complete nitrate removal at low flow (2.9 cm d$^{-1}$) declining to 30\% at 13.6 cm d$^{-1}$ \parencite{Greenan2009}, and 100\% removal over 100 days at 20 mg N L$^{-1}$ and 24-h HRT (10.1 g N m$^{-3}$ d$^{-1}$) \parencite{Hua2016}.

\subsection{Environmental and hydraulic controls on performance}

The dominant, repeatedly quantified control is hydraulic residence time, which trades off against the fraction of flow captured. Removal efficiency rises with HRT across studies---from 20\% (2 h) to 98\% (8 h) in Illinois \parencite{Bell2015}, from 31\% to 93\% over 2.8--6.4 d in New Zealand \parencite{Rivas2020}, and with efficiency increasing at longer daily HRT but trending negative below about 2 h in a paired Illinois system \parencite{Maxwell2022}. In the ten-bioreactor Illinois dataset, annual flow-weighted HRT correlated strongly with flow captured ($\rho=-0.67$) and with treated-water N removal efficiency ($\rho=0.60$) but not with edge-of-field efficiency ($\rho=0.03$), and increased HRT increased efficiency for treated water as the standardised-design theory predicts \parencite{Christianson2026}. Column work shows removal rate per unit wood increasing with flow while efficiency falls ($y=4.299x^{-1}$, $r^2=0.96$) \parencite{Greenan2009}, and heated-treatment tests showed removal rising from 1.5 mg NO$_3$-N L$^{-1}$ at 1 h to $>$18.6 mg L$^{-1}$ at 8 h \parencite{Ma2021}. Influent load also matters: higher loads raise removal \emph{rate} but lower percentage efficiency and outlet-concentration reductions, a pattern seen for bioreactors \parencite{Rivas2020,Christianson2026} and echoed in constructed-wetland syntheses \parencite{Mendes2021}. A meta-analysis of solid-carbon denitrification similarly found higher influent nitrogen associated with higher rates, and---for those studies---inferior removal above 5-h HRT \parencite{Zhang2021b}.

Temperature exerts a strong, consistent effect. In the only controlled Illinois field experiment holding influent nitrate and HRT constant, exponential functions of water temperature and HRT explained 73\% of the variance in load reduction and 43\% in removal rate \parencite{Bell2015}. In South Dakota, average removal was 56\% in March--June versus 81\% in July--August, falling by about 30\% below 12\,$^{\circ}$C \parencite{Thapa2023}; laboratory columns reached 96\% at 21.5\,$^{\circ}$C/24-h HRT versus 29\% at 10\,$^{\circ}$C/12-h HRT \parencite{Soupir2018}; and Michaelis--Menten $V_{max}$ dropped sharply from 22 to 5\,$^{\circ}$C for composted and fresh woodchips \parencite{Kouanda2021}. Cross-study analysis found low temperature ($<$20\,$^{\circ}$C) significantly depressed removal while intermediate and high temperatures did not differ \parencite{Zhang2021b}. Flow dynamics and antecedent conditions further modulate performance: removal was lower during high-flow events and at high inflow nitrate, and samples taken immediately after rainfall showed higher load removal than later samples \parencite{Thapa2023}. In Quebec, influent nitrate correlated with soil moisture ($r=0.64$) and prior-day temperature ($r=0.67$) and negatively with flow rate ($r=-0.53$), while in-bioreactor reduction tracked lagged temperature ($r=0.58$) and influent concentration ($r=0.71$), with no clear HRT--percent-removal relationship at that site \parencite{Husk2017}. Reaction kinetics underpin these observations: removal follows zero-order kinetics when nitrate is non-limiting (rate constants 1.42--1.80 mg N L$^{-1}$ h$^{-1}$) and shifts to first-order ($0.49$ h$^{-1}$) once nitrate becomes limiting \parencite{Hua2016}, and can be treated operationally as zero-order when nitrate greatly exceeds $K_m$ \parencite{Kouanda2021}.

\subsection{Carbon substrates, media, and removal mechanisms}

Denitrification is firmly established as the dominant nitrate-removal pathway. In Iowa microcosms it accounted for $>$98\% of transformed NO$_3$-N in all treatments (except cardboard, $>$96\%), with immobilisation $<$2.4\% and DNRA $<$1\% \parencite{Greenan2006}, and flow-through columns confirmed denitrification as the primary fate with immobilisation minor \parencite{Greenan2009}. Denitrification potential in woodchips is sustained and far exceeds surrounding soil: over nine years wood increased denitrification potential 31- to 4000-fold over subsoil, denitrifier populations exceeded $10^8$ g$^{-1}$ wood, and in situ push--pull testing gave nitrate half-lives of 0.64--0.81 d, implying 24-h HRT suffices to bring 20--25 mg L$^{-1}$ influent below 10 mg NO$_3$-N L$^{-1}$ \parencite{Moorman2010}. Reviews likewise identify denitrification as the main mechanism ahead of biological uptake and settling \parencite{Mendes2021}, while noting that DNRA (favoured under low NO$_3$-N and high carbon) and, more tentatively, anammox also occur \parencite{Hartfiel2022}.

Carbon quality and longevity govern long-term performance. Cellulose is the most rapidly consumed woodchip compound, making cellulose metabolism central to persistence \parencite{Hartfiel2022}. Woodchip decomposition is depth- and position-dependent: in an Iowa wall, wood loss reached 75\% at 90--100 cm (half-life 4.6 y) but $<$13\% at 155--170 cm (half-life 36.6 y) \parencite{Moorman2010}, and along a Minnesota bed the inlet showed greater decomposition (lower C/N, higher lignocellulose index) than downstream, with aged four-year woodchips more recalcitrant than fresh \parencite{Ghane2018}. Alternative and blended substrates can raise performance. Corncob-amended systems removed more nitrogen than woodchips alone: pilot CC75 reached 61.2\% at 16-h HRT versus 37.9\% for WC100, and modelled full-scale CC75 gave the highest annual load removal (15.6\%) and lowest cost, remaining cost-advantageous even with periodic corncob replacement \parencite{Law2023}. Corn stover blended up to 10\% with woodchips enhanced denitrification (mean concentration reductions of 86\% for corn stover, 89\% for the blend, 77\% for woodchips), though pure corn stover raised ammonium and total inorganic N and is not recommended as the main medium \parencite{Oladeji2023}. Microcosm rankings placed cornstalks $>$ cardboard $>$ wood plus soybean oil $>$ woodchips alone, with grinding and soybean-oil additions boosting removal \parencite{Greenan2006}, and Canadian-prairie testing found agricultural residues (barley straw best) outperforming woodchips for nitrate removal, though woodchips gave better hydraulic efficiency \parencite{Kohn2023}. Amendments further tune performance: activated carbon raised removal at low temperature ($<$10\,$^{\circ}$C) and cut organic-carbon losses while flaxseed cake did not improve removal despite greater carbon release \parencite{Povilaitis2020}; biochar increased nitrate removal in proportion to dose \parencite{Ahmadvand2020}; and composting woodchips before installation raised denitrification capacity 1.20--2.15-fold at 22\,$^{\circ}$C with better cold tolerance \parencite{Kouanda2021}. Broadly, synthetic solid carbon sources outperform natural ones, while mixed and single carbon sources did not differ significantly in a meta-analysis, indicating species matters more than mixing \parencite{Zhang2021b}. A recurring minor by-product is nitrite accumulation (9.6--18.7\% of reduced nitrate), which increases as HRT shortens and influent nitrate rises \parencite{Hua2016,Kouanda2021}.

\subsection{Microbial ecology of denitrifying bioreactors}

Bioreactor microbial communities are variable and shaped by substrate availability, depth, season, moisture, pH, temperature, inlet nitrate, dissolved oxygen, and HRT \parencite{Hartfiel2022}. Culture and molecular work has identified the responsible organisms, especially those active in cold conditions. From Minnesota bioreactors, screening of 206 isolates yielded three confirmed denitrifiers (\textit{Cellulomonas} sp. WB94, \textit{Delftia} sp. BE1.2, \textit{Microvirgula} sp. BE2.4), two of which denitrified aerobically and whose genomes carried denitrification and polysaccharide-degradation genes, marking WB94 and BE2.4 as bioaugmentation candidates \parencite{Anderson2020}. A parallel study identified denitrifiers active at cold woodchip conditions---enriching \textit{Pseudomonas}, \textit{Polaromonas}, and \textit{Cellulomonas}---and confirmed \textit{Cellulomonas} sp. WB94 (98\% ANI to \textit{C.\ cellasea}) reduces nitrate and degrades cellulose \parencite{Jang2019}. Nitrate-reducing fungi are also abundant and cold-adapted: densities reached 845,000 and 524,000 CFU g$^{-1}$ in soil and woodchips at 15\,$^{\circ}$C, 91 of 349 isolates produced gaseous N (frequently $^{29}$N$_2$, suggesting co-denitrification), 89\% showed cellulase activity, and \textit{nirK} and \textit{p450nor} were detected in 38\% and 12\% of strains \parencite{Aldossari2021}. Bioaugmentation shows conditional promise: pre-grown \textit{Pseudomonas} isolates from Danish woodchips did not help in the first assay but increased nitrate removal by 27--37\% at 5\,$^{\circ}$C in the second, with \textit{P.\ chlororaphis} PSS4 most effective \parencite{Jeglot2022}; however, inoculation alone often fails owing to rapid washout \parencite{Wang2025b}.

Community function tends to be robust. Under fluctuating water levels, denitrification rates and denitrifier communities were generally resistant to disturbance, and nitrous-oxide production rates were consistently about ten-fold lower than nitrate-removal rates \parencite{Hathaway2017}. Mineral amendment can steer the community: heated tourmaline reduced biofilm growth (lowest DNA) yet increased the relative abundance of \textit{nirS}-, \textit{nirK}-, and \textit{nosZ}-encoding denitrifiers, reducing clogging while giving the highest removal rates (52.1, 43.1, 48.0 g NO$_3$-N m$^{-3}$ d$^{-1}$ at 1, 4, 8 h) \parencite{Ma2021}. Biofouling deposits in drainage wells were controlled chiefly by dissolved iron rather than aerobic bacterial presence or dissolved oxygen, with stalked bacteria only where fouling occurred \parencite{Stuetz2004}.

\subsection{Bioreactor design, sizing, and hydraulics}

Sizing procedures aim to balance flow treatment against removal efficiency. Using 20\% of the Manning's-estimated peak flow generally treated the majority of annual volume (54\% at Hamilton, 72\% at SERF), and modelling supports design guidance of 10--20\% of peak flow at 6--8-h design retention times; for fixed volume, longer, narrower designs removed a greater percentage of load and had higher rates than short, wide designs \parencite{Christianson2013}. Load-duration analysis across three watersheds---where roughly 50\% of annual NO$_3$-N moved in the largest 30\% of daily base flows---showed that bioreactors sized to a 0.5-d HRT could remove about 44--52\% of base-flow nitrate and at least 20--30\% of total annual loss while occupying at most 0.27\% of watershed area \parencite{Moorman2015}. Notably, treating the majority of annual flow is not universally necessary where goals are load-based: a similar design flow rate ($\sim$1.7--1.8 L s$^{-1}$) corresponded to 50\% non-exceedance in small drainage areas but only 5\% in large ones, so treated volumes differed greatly by area \parencite{Christianson2025}. Loading-density guidance from the Conservation Practice Standard was not supported by the ten-bioreactor data, which showed no significant relationship between loading density and removal efficiency \parencite{Christianson2026}.

Hydraulic behaviour links design to performance. Across 14 tracer tests at six full-size bioreactors, high-aspect-ratio conventional designs behaved most like plug flow (tanks-in-series), the wide design showed the most dispersion and short-circuiting, and an advanced wide design with baffles achieved higher volumetric efficiency and trended toward the highest removals; dispersion/short-circuiting and volumetric/hydraulic-efficiency metrics were each strongly correlated \parencite{deOliveira2023}. A wide baffled Illinois bioreactor nonetheless did not raise the fraction of flow treated ($\sim$40\%) or removal relative to conventional designs---attributed partly to flashy drainage (Richards--Baker index 0.74--0.91)---yet still achieved consistent 22--24\% edge-of-field reduction with 54--69.5\% internal efficiency and clear along-path nitrate declines \parencite{Dougherty2020b}. Expanding cross-sectional area improved flow capture: a paired short-and-wide system captured 84--92\% of annual discharge from a 29-ha field \parencite{Maxwell2022}. Wood-filled bioreactors showed better hydraulic efficiency than other feedstocks in the Canadian prairies \parencite{Kohn2023}. For buffer-diverting designs, more diverted flow did not mean more nitrate removal, buffer width reduced diversion nonlinearly, and neglecting exit head loss overestimated diverted flow and removal, motivating a process-based design that includes a removal component \parencite{Abdalaal2023}. Physical durability is quantifiable: hand-held LiDAR at 17 full-size bioreactors captured surface subsidence, which correlated with age (median 1.0 cm y$^{-1}$ at soil-covered versus 7.3 cm y$^{-1}$ at non-covered sites; maximum 73 cm) and with proximity to the drainage outlet \parencite{Christianson2025b}.

\subsection{Design innovations and performance-enhancement strategies}

Several modifications extend or intensify treatment. Pumping can supply additional nitrate-laden water and lengthen the operating season: a retrofitted dual-chamber system pumped at 1.06--2.17 L s$^{-1}$ per chamber, extended operation by three weeks, and raised the combined daily removal rate to 6.58 g N m$^{-3}$ d$^{-1}$ (from 3.62 without pumping), albeit bypassing 41\% of pumped flow untreated \parencite{Hartfiel2023}. At an underloaded Illinois bioreactor, pumping increased treated volume, loading, and removal---removing 2.8 kg N in a 29-d 2022 period (versus 0.69--2.0 kg N unpumped) and 7.2 kg N over 93 d in 2023---at daily rates of 7.5 and 5.2 g N m$^{-3}$ d$^{-1}$ and efficiencies of 57\% and 64\%, with nightly batching giving 100\% instantaneous removal and a small heating benefit \parencite{DíazGarcía2025}. Electrical stimulation improved denitrification at 500 mA (carbon anode--cathode best) but not at 100 mA, increased efficiency at the cost of current efficiency, and raised removal cost, so gains in efficiency would be needed for cost-effectiveness \parencite{Law2018}. Heated tourmaline both reduced clogging and gave the highest removal across tested HRTs \parencite{Ma2021}. Scaling up, a three-bed catchment-scale Minnesota bioreactor treated 55\% of catchment discharge and removed 713 kg nitrate-N (12.5\% of catchment load; 23\% of throughflow load), but sediment clogging during high flow reduced performance and required a sediment-exclusion system \parencite{Feyereisen2023}. Beyond tile drainage, spring bioreactors in the Chesapeake region averaged 41\% annual load reduction (up to $>$65\% in summer), and treating hundreds of springs was estimated to remove 322--771 kg d$^{-1}$ at low unit cost \parencite{Easton2019}. Non-biological alternatives include anion-exchange resin, which removed nitrate at both low and high flow with near-identical breakthrough volumes (favouring high-flow scale-up) but was hindered by sulfate selectivity and regenerant-disposal issues \parencite{MAGETTE1990}. Reviews add that carbon dosing and woodchip--corncob mixtures can substantially raise performance, though carbon dosing risks leaving unused effluent carbon and causing bioclogging or pollution swapping \parencite{Wang2025b}.

\subsection{Pollution swapping and greenhouse-gas by-products}

Nonideal conditions can generate N$_2$O, CH$_4$, CO$_2$, sulfate reduction, and methylmercury, but reported risks are generally low \parencite{Hartfiel2022}. For nitrous oxide, column and field studies consistently find emissions small relative to nitrate removed: N$_2$O production was 0.003--0.028\% of N denitrified in columns \parencite{Greenan2009}; an Iowa wall showed no significant N$_2$O difference from controls, indirect emission factors of $\sim$0.0085--0.0088 kg N$_2$O-N kg$^{-1}$ NO$_3$-N, and indirect losses equivalent to tile drainage without the bioreactor \parencite{Moorman2010}; and a Virginia bioreactor had an average N$_2$O emission factor of 0.07\% (median 0.01\%), well below the 0.75\% IPCC comparison, with GHG emissions not environmentally concerning \parencite{Bock2018}. Across five Illinois bioreactors, dissolved N$_2$O was produced at four sites but consumed at one, surface N$_2$O fluxes were much smaller than dissolved fluxes, only a small fraction of removed nitrate was emitted as N$_2$O, and modelled watershed-scale indirect N$_2$O emissions were lower with a bioreactor than without \parencite{Brunton2024}. Methane and sulfate responses are more variable. In Illinois, three of five bioreactors produced notable CH$_4$ (which drove global-warming potential where present) while soil covers significantly lowered surface CH$_4$ \parencite{Brunton2024}; in an Iowa dual-chamber system both chambers were significant CH$_4$ sources, one chamber removed sulfate and dissolved N$_2$O while the other did not, yet the landscape-scale impact was a $<$1\% increase in CO$_2$e, with pollution swapping more likely above 50-h HRT \parencite{Hartfiel2023b}. Sulfate reduction under non-N-limiting conditions is common and increases with HRT---occurring in 83\% of Illinois runs \parencite{Bell2015}, highest in activated-carbon media and associated with lower nitrate removal at long HRT \parencite{Povilaitis2020}, and indicating strongly reduced, potentially H$_2$S-producing conditions in New Zealand \parencite{Rivas2020}. Carbon leaching, especially at start-up, is a further concern: DOC shows an initial flush then decline \parencite{Abusallout2017,Rivas2020}, is dominated by hydrophobic, high-molecular-weight (10--100 kDa) fractions that are strong disinfection-by-product precursors with TOX yields comparable to Suwannee River Fulvic Acid, and can be removed by steel-chip filters or coagulation (alum and PAC more effective than ferric chloride) \parencite{Abusallout2017}. Ammonium and DOC also leach at higher flow \parencite{Greenan2009}, and a catchment-scale system net-exported 212 kg C (22\% of inlet DOC load) \parencite{Feyereisen2023}.

\subsection{Phosphorus co-removal and enhanced media}

Because woodchips remove little phosphorus---and often leach it at start-up---phosphorus co-treatment has driven the development of coupled media \parencite{Christianson2021}. Steel-based media are consistently effective P sorbents. Steel turnings achieved $>$90\% orthophosphate removal, and placing iron both upstream and downstream of woodchips overcame arrangement conflicts to enhance both nitrate and orthophosphate removal (woodchips reduce nitrate best downstream of iron, while iron re-captures P released by woodchips) \parencite{Goodwin2015}. Steel byproducts showed batch Langmuir $q_{max}$ of 1.30--5.81 mg P g$^{-1}$ and a continuous-flow capacity of 3.70 mg P g$^{-1}$, completely removing residual phosphate (1.0 g P m$^{-3}$ d$^{-1}$) and reducing effluent nitrite in a two-stage woodchip-then-steel system \parencite{Hua2016}. Fly-ash pellets removed on average 71.4\% of orthophosphate while paired woodchip--pellet configurations held total phosphate removal at 81.7--84.5\% \parencite{Li2018b}, and iron filings cut drainage P loads by 19\% alongside nitrate removal \parencite{Oladeji2023}. Field bioreactors can also become net P sinks over time: Quebec woodchips-only and mixed-media systems released P initially but became sinks within a year, with mixed media removing $\sim$19 times more total-P load and cutting soluble reactive P by 97\% (versus 53\% for woodchips only), though neither met the 0.030 mg L$^{-1}$ criterion \parencite{Husk2018}. A New Zealand bioreactor unexpectedly retained DRP in its second year \parencite{Rivas2020}, and a catchment-scale system removed DRP and TP for most events but exported TP at peak flow \parencite{Feyereisen2023}. Some systems are inconsistent: a Minnesota modular design acted as both P sink and source depending on month and media (fragmented limestone $>$ crushed concrete $>$ steel slag for apparent P-sorbing capacity), with media requiring annual replacement \parencite{Strock2025}, and P removal was not statistically significant in the Virginia system \parencite{Bock2018}. A promising biological route is drying--rewetting: cyclic draining and re-flooding of woodchip columns more than doubled orthophosphate removal (peaking 10--12 h after re-flooding), driven by biological uptake---confirmed because autoclaved woodchips did not take up phosphate and \textit{ppk} transcripts of \textit{Accumulibacter} and \textit{Pseudomonas} peaked with uptake---offering a simple co-treatment strategy for N and P \parencite{Zhang2024b}. Microbial P uptake, rather than woodchip sorption alone, may also explain variable DRP behaviour in columns \parencite{Soupir2018}.

\subsection{Secondary water-quality benefits and monitoring context}

Beyond nitrate, woodchip bioreactors provide co-benefits. Columns effectively removed \textit{Salmonella} and \textit{E.\ coli} under all conditions (greater at 21.5 than 10\,$^{\circ}$C), with temperature more influential than HRT, indicating a pathogen-removal co-benefit for systems run for nitrate \parencite{Soupir2018}. Reliable performance assessment depends on well-defined monitoring: bioreactors offer discrete influent and effluent points for direct measurement \parencite{Easton2019}, and standardised monitoring set-ups and minimum reporting measures have been recommended to compare annual load-removal performance \parencite{Christianson2021}. Supporting analytical methods have also been validated---for example the Dumas method for total N and $^{15}$N ratios of finely ground plant material and evaporatively concentrated drainage samples (highest reported recovery 97.4\%, improved with copper shot) \parencite{Schindler1999}---and resin-capsule methods can intercept nitrate moving via macropore and bypass flow, the primary transport pathway to subsurface drainage \parencite{Pampolino2000}.

\subsection{Cost-effectiveness and comparison with other edge-of-field practices}

Bioreactors are generally low-cost per unit nitrogen removed. Installation costs have ranged from under \$5,000 to \$27,000, with cost efficiencies generally from $<$\$2.50 to $\sim$\$20 kg$^{-1}$ N y$^{-1}$ (up to \$48) \parencite{Christianson2021}. Corncob (CC75) systems were the most cost-efficient carbon option (\$10.56--13.89 kg$^{-1}$ N, versus up to \$88.11 for woodchips) and remained advantageous even under conservative replacement assumptions \parencite{Law2023}. For pumped systems, unit cost fell as lifespan, mass-removal rate, and operating period rose---raising mass-removal rate from low to high cut unit cost by 91\%---with best-case costs of $\sim$\$4.80--23.81 kg$^{-1}$ NO$_3$-N (worst case $\sim$\$140--619), and increasing removal rate and operating period beyond six months identified as the strongest cost levers \parencite{Hartfiel2024}. Spring bioreactors achieved \$0.54--7.60 kg$^{-1}$ N y$^{-1}$, typically $<$\$5 for higher-flow, higher-concentration springs \parencite{Easton2019}, whereas electrical stimulation raised costs above the base case \parencite{Law2018}.

Set against other edge-of-field practices, bioreactors are effective for nitrate but comparatively limited in co-benefits. A US Corn Belt review of 119 records found all five practices (controlled drainage, water-quality wetlands, bioreactors, drainage water recycling, saturated buffers) effective at removing nitrate, but with differing additional benefits and disservices: water-quality wetlands gave the highest nitrate reduction plus phosphorus removal and habitat value, drainage water recycling offered the greatest yield and co-benefit potential, and denitrifying bioreactors removed nitrate effectively but with limited extra benefits and a need for further study of phosphorus, pesticides, and greenhouse gases \parencite{Mitchell2023}. Global synthesis places bioreactor flow-season removal at 14--66\%, alongside constructed wetlands ($\sim$39\% average TN removal) and saturated buffers and controlled drainage ($\sim$30--50\% annual N-load reduction), while noting that agricultural drainage's high-nitrate, low-carbon, low-temperature character complicates direct transfer of municipal treatment approaches \parencite{Wang2025b}. Constructed-wetland and paired-system studies reinforce the theme that design and sizing dominate outcomes: about 40--50\% N removal is generally achievable at area ratios of 1--4\% with minor gains beyond \parencite{Mendes2021}, and paired surface-flow wetland plus woodchip systems acted as nutrient sinks (TN removal 8--51\%, TP 0--80\%), with the woodchip component contributing more removal than the sedimentation basin and requiring smaller footprints than standalone wetlands \parencite{Pugliese2020}. Collectively, this evidence establishes woodchip bioreactors as a dependable, cost-effective nitrate-reduction practice whose performance is well predicted by HRT, temperature, flow, and loading, and whose principal remaining design frontiers concern phosphorus, greenhouse gases, and pesticides.